\documentclass[tikz,border=3mm]{standalone}

%%%% Serifenloser Textstil (Calibri-ähnlich)
\usepackage[T1]{fontenc}
\usepackage[sfdefault,lf]{carlito}
\usepackage{lmodern}
\renewcommand*\familydefault{\sfdefault}

\usepackage{graphicx}
\usepackage{xcolor}
\usetikzlibrary{positioning,calc,arrows.meta,shapes.geometric,fit}

% =========================================================
% ======= Regler / Parameter (oben zentral anpassen) ======
% =========================================================
% --- Flow-Boxen ---
\def\flowW{5.0}        % Breite der Haupt-Flowboxen [cm]
\def\flowH{1.05}       % Mindesthöhe der Haupt-Flowboxen [cm]
\def\subW{4.0}         % Breite der Plan-B-Substeps [cm]

\def\titleSep{0.65}    % Abstand Titel -> erste Box [cm]

% Vertikale Abstände (Row-Gaps)
\def\mainVsep{1.2}     % Abstand nach einer "Row" [cm]
\def\subVsep{1.0}      % Abstand zwischen Plan-B-Substeps [cm]

% --- Feinjustage ---
\def\decSep{0.25}       % Abstand PlanA-ROW -> Decision [cm] (kleiner)
\def\rotUp{2.0}         % [cm] Rotations-Box nach oben (positiv = hoch)
\def\planBLabelX{1.6}  % [cm] Label "Plan B (Fallback)" nach rechts

% --- Bildspalte rechts ---
\def\imgW{3.5}         % Bildbreite rechts [cm] (Höhe ist dynamisch!)
\def\imgGap{1.2}       % Horizontaler Abstand Flow -> Bildspalte [cm]

% Placeholder-Höhe (nur falls Datei fehlt)
\def\phH{2.2}          % [cm]

% --- Bypass (ja-Pfad) links am Plan-B-Kasten vorbei ---
\def\leftBypassPad{0.7} % wie weit links vom Plan-B-Groupbox der ja-Pfad laufen soll [cm]

% --- Plan-B-Block verschieben (kompletter Fallback-Block) ---
\def\planBX{0.0}       % [cm] + rechts / - links
\def\planBY{0.0}       % [cm] + hoch / - runter

% --- Linien ---
\def\boxLW{0.8pt}
\def\flowLW{0.85pt}
\def\imgFrameLW{0.65pt}
\def\imgLinkLW{0.65pt}
\def\groupLW{0.65pt}
\def\boxInner{4pt}

% =========================================================
% ======= Bildpfade (hier einfach ersetzen) ===============
% =========================================================
\def\imgCrop{Figures/step_crop.png}
\def\imgPlanA{Figures/step_planA.png}
\def\imgPlanBFilters{Figures/step_PlanB_Filters.png} % vom Groupbox verlinkt
\def\imgPlanB{Figures/step_planB.png}                % danach (lila Box)
\def\imgWarp{Figures/step_warp.png}

% =========================================================
% ======= Bild-Offsets je Schritt (X/Y in cm) =============
% =========================================================
\def\dxCrop{0.0}    \def\dyCrop{1.0}
\def\dxPlanA{0.0}   \def\dyPlanA{-0.2}
\def\dxFilters{0.3} \def\dyFilters{1.3}
\def\dxPlanB{0.5}   \def\dyPlanB{-0.4}
\def\dxWarp{0.0}    \def\dyWarp{-1.0}

% =========================================================
% ======= Helper: Bild oder Platzhalter ===================
% =========================================================
\newcommand{\ImgOrPlaceholder}[1]{%
  \begingroup
  \def\file{#1}%
  \IfFileExists{\file}{%
    \includegraphics[width=\imgW cm]{\file}%
  }{%
    \begin{minipage}[c][\phH cm][c]{\imgW cm}\centering
      \color{black!55}\scriptsize
      Bild-Platzhalter\\[2pt]
      \texttt{\file}
    \end{minipage}%
  }%
  \endgroup
}

% Step -> Bild rechts + gestrichelte Verbindung (mit X/Y-Shift)
\newcommand{\PlaceImageShift}[5]{% (stepnode, imagenode, file, dx, dy) in cm
  \node[imgframe, anchor=west] (#2)
    at ($([xshift=\imgGap cm]#1.east)+(#4 cm,#5 cm)$) {\ImgOrPlaceholder{#3}};
  \draw[imglink] (#1.east) -- (#2.west);
}

% Row-Bottom: nimmt das "unterste" von Box+Bild (min y) als Referenz
\newcommand{\RowSouthTwo}[3]{% (nodeA, nodeB, coordname)
  \path let \p1=(#1.south), \p2=(#2.south) in
    coordinate (#3) at (\x1,{min(\y1,\y2)});
}

% Row-Bottom aus drei Knoten
\newcommand{\RowSouthThree}[4]{% (nodeA, nodeB, nodeC, coordname)
  \path let \p1=(#1.south), \p2=(#2.south), \p3=(#3.south) in
    coordinate (#4) at (\x1,{min(\y1,\y2,\y3)});
}

\begin{document}
\begin{tikzpicture}[
  >={Latex[length=2.8mm]},
  every node/.style={font=\scriptsize},
  flow/.style={-Latex, line width=\flowLW},
  imglink/.style={dashed, line width=\imgLinkLW, color=black!70},
  flowbox/.style={
    draw=black, rounded corners=3pt, line width=\boxLW,
    align=left, fill=white,
    minimum width=\flowW cm, minimum height=\flowH cm,
    inner sep=\boxInner
  },
  subbox/.style={
    flowbox, minimum width=\subW cm, rounded corners=2pt
  },
  decision/.style={
    draw=black, diamond, aspect=2.0,
    inner sep=1.5pt, line width=\boxLW, fill=white, align=center
  },
  imgframe/.style={
    draw=black, line width=\imgFrameLW, inner sep=0pt,
    minimum width=\imgW cm, fill=white
  },
  groupbox/.style={
    draw=black!65, dashed, rounded corners=4pt, line width=\groupLW,
    inner sep=6pt
  },
  edgelabel/.style={font=\scriptsize, fill=white, inner sep=1.2pt}
]

% ---------- Titel ----------
\coordinate (L0) at (0,0);
\node[font=\bfseries\small, anchor=north west] (title) at (L0)
  {Pipeline-Überblick (Plan A/B) und Zielgrößen};

% =========================================================
% ====================== FLOW LINKS =======================
% =========================================================
\coordinate (cursorNW) at ($(title.south west)+(0,-\titleSep cm)$);

\node[flowbox, anchor=north west] (in) at (cursorNW) {%
  \textbf{Eingabe:} Vollbild $\rightarrow$ Kandidatenbox\\
  (ROI/CNN-Merging)
};

% nächste Zeile unter "in"
\coordinate (cursorNW) at ($(in.south west)+(0,-\mainVsep cm)$);

\node[flowbox, anchor=north west] (crop) at (cursorNW) {%
  \textbf{Crop:} Kandidatenbereich\\
  mit kleinem Padding ausschneiden
};
\PlaceImageShift{crop}{cropImg}{\imgCrop}{\dxCrop}{\dyCrop}
\RowSouthTwo{crop}{cropImg}{cropRowSouth}

% nächste Zeile unter Crop-ROW
\coordinate (cursorNW) at ($(in.west |- cropRowSouth)+(0,-\mainVsep cm)$);

\node[flowbox, anchor=north west] (planA) at (cursorNW) {%
  \textbf{Plan A:} Direkte Eckpunktdetektion\\
  \texttt{cv2.QRCodeDetector.detect()}
};
\PlaceImageShift{planA}{planAImg}{\imgPlanA}{\dxPlanA}{\dyPlanA}
\RowSouthTwo{planA}{planAImg}{planARowSouth}

% Decision: unter PlanA-ROW (zentriert)
\node[decision, anchor=north] (ok) at ($(planARowSouth)+(0,-\decSep cm)$) {%
  Plan A\\
  erfolgreich?
};

% =========================================================
% =================== PLAN B (SUBSTEPS) ===================
%   -> exakt unter Decision (gerader nein-Pfeil)
%   -> optional verschiebbar via \planBX/\planBY
% =========================================================
\coordinate (bTop) at ($(ok.south)+(0,-\subVsep cm)+(\planBX cm,\planBY cm)$);

\node[subbox, anchor=north] (b1) at (bTop) {%
  \textbf{Plan B:} Binarisierung
};

\node[subbox, anchor=north] (b2) at ($(b1.south)+(0,-\subVsep cm)$) {%
  Morphologie
};

\node[subbox, anchor=north] (b3) at ($(b2.south)+(0,-\subVsep cm)$) {%
  Kontur / Hull
};

\node[subbox, anchor=north] (b4) at ($(b3.south)+(0,-\subVsep cm)$) {%
  Viereck via \texttt{approxPolyDP}\\
  (mit IoU-Scoring)
};

\node[groupbox, fit=(b1)(b4),
      label={[font=\scriptsize\bfseries, xshift=\planBLabelX cm]north west:Plan B (Fallback)}] (bGroup) {};

% Bild für Plan-B-Filters: vom GROUPBOX aus
\PlaceImageShift{bGroup}{bFiltersImg}{\imgPlanBFilters}{\dxFilters}{\dyFilters}

% Bild für Step Plan B (lila Box): vom Ende (b4) aus
\PlaceImageShift{b4}{bPlanBImg}{\imgPlanB}{\dxPlanB}{\dyPlanB}

% Row-Bottom Plan-B-Bereich (Groupbox + beide Bilder)
\RowSouthThree{bGroup}{bFiltersImg}{bPlanBImg}{bRowSouth}

% =========================================================
% ===================== WARP (+ ROT) ======================
% =========================================================
\coordinate (warpNW) at ($(in.west |- bRowSouth)+(0,-\mainVsep cm)$);

\node[flowbox, anchor=north west] (warp) at (warpNW) {%
  \textbf{Warp:} Perspektivische Entzerrung\\
  des Vierecks $\rightarrow$ quadratische Top-View
};
\PlaceImageShift{warp}{warpImg}{\imgWarp}{\dxWarp}{\dyWarp}
\RowSouthTwo{warp}{warpImg}{warpRowSouth}

% Rotationen (ohne Bild) - nach oben verschiebbar via \rotUp
\coordinate (rotNW) at ($(in.west |- warpRowSouth)+(0,-\mainVsep cm)+(0,\rotUp cm)$);
\node[flowbox, anchor=north west] (rot) at (rotNW) {%
  \textbf{Rotationen:} Schätze $(rot_x, rot_y, rot_z)$\\
  aus dem Viereck
};

% =========================================================
% ====================== PFEILE ===========================
% =========================================================
\draw[flow] (in.south) -- (crop.north);
\draw[flow] (crop.south) -- (planA.north);
\draw[flow] (planA.south) -- (ok.north);

% nein: gerade nach unten (weil b1 unter ok zentriert ist)
\draw[flow] (ok.south) -- node[edgelabel, left] {nein} (b1.north);

% Plan B chain
\draw[flow] (b1.south) -- (b2.north);
\draw[flow] (b2.south) -- (b3.north);
\draw[flow] (b3.south) -- (b4.north);

% ja: links am Plan-B-Groupbox vorbei und auf warp.west rein
\coordinate (leftLaneX) at ($(bGroup.west)+(-\leftBypassPad cm,0)$);
\coordinate (laneTop)   at (leftLaneX |- ok.west);
\draw[flow]
  (ok.west) -- node[edgelabel, near start, above] {ja} (laneTop)
  |- (warp.west);

% Weiter nach unten (Plan B -> Warp)
\draw[flow] (b4.south) -- (warp.north);
\draw[flow] (warp.south) -- (rot.north);

\end{tikzpicture}
\end{document}
